% ==================================================================
% CONFIGURAÇÕES INSTITUCIONAIS E PACOTES
% Este arquivo consolida todas as configurações visuais, pacotes e 
% macros do Instituto Federal de Sergipe (IFS).
% ==================================================================

% --- Codificação e Tipografia Moderna ---
\usepackage{iftex}
\ifPDFTeX
    \usepackage[utf8]{inputenc}
    \usepackage[T1]{fontenc}
    \pdfminorversion=7 % Permite incluir figuras PDF 1.7 (Python, Mermaid, Inkscape, Canva) sem avisos
\fi

% --- Fonte acadêmica moderna: Palatino ---
\usepackage{mathpazo} % Tipografia Palatino clássica (inclusa no TeX Live base)
\usepackage[
    activate={true,nocompatibility},
    tracking=true,
    factor=1000,
    stretch=10,
    shrink=10
]{microtype}
\microtypecontext{spacing=nonfrench}

\usepackage{indentfirst}
\usepackage{color}
\usepackage{graphicx}
\graphicspath{{figuras/}{figuras/diagramas/}{figuras/graficos/}{figuras/capturas/}}
\usepackage{pdfpages} % Permite incluir o PDF da Ficha Catalográfica emitida pela biblioteca
\usepackage{csquotes}

% --- Matemática e Tabelas ---
\usepackage{amsmath}
\usepackage{booktabs}
\usepackage{tabularx}
\usepackage{multirow}
\usepackage{array}
\usepackage{longtable}
\usepackage{siunitx}
\usepackage[section,below]{placeins}
\usepackage{etoolbox}
\pretocmd{\subsection}{\FloatBarrier}{}{} % Impede flutuadores de vazarem para a subseção seguinte

% --- Parâmetros de Diagramação de Flutuadores (Tabelas e Figuras) ---
\renewcommand{\topfraction}{0.85}
\renewcommand{\bottomfraction}{0.70}
\renewcommand{\textfraction}{0.15}
\renewcommand{\floatpagefraction}{0.70}
\setcounter{topnumber}{3}
\setcounter{bottomnumber}{2}
\setcounter{totalnumber}{5}

% Configuração nativa do memoir para legendas ABNT (evita usar o pacote caption que conflita)
\captionnamefont{\small\bfseries}
\captiontitlefont{\small}
\captiondelim{ -- } % Padrão ABNT (travessão separando label do texto)

% --- Ambiente Formal de Quadros (ABNT / IBGE - 100% nativo do memoir/abnTeX2) ---
\newcommand{\quadroname}{Quadro}
\newcommand{\listofquadrosname}{LISTA DE QUADROS}
\newfloat[chapter]{quadro}{loq}{\quadroname}
\newlistof{listofquadros}{loq}{\listofquadrosname}
\newlistentry{quadro}{loq}{0}
\counterwithout{quadro}{chapter}
\renewcommand{\thequadro}{\arabic{quadro}}
\renewcommand{\cftquadroname}{\quadroname\space}
\renewcommand*{\cftquadroaftersnum}{\hfill--\hfill}
\providecommand*{\quadroautorefname}{Quadro}

% --- Código-fonte ---
\usepackage{listings}
\usepackage{xcolor}
\usepackage{courier} % Fonte monospace padrão de alta compatibilidade
\renewcommand{\lstlistingname}{Código}
\definecolor{codeBackground}{HTML}{F8F9FA}
\definecolor{codeFrame}{HTML}{DEE2E6}
\definecolor{codeKeyword}{HTML}{0550AE}
\definecolor{codeString}{HTML}{0A3069}
\definecolor{codeComment}{HTML}{6E7781}
\definecolor{codeNumber}{HTML}{953800}
\lstset{
    backgroundcolor  = \color{codeBackground},
    basicstyle       = \ttfamily\small,
    keywordstyle     = \color{codeKeyword}\bfseries,
    stringstyle      = \color{codeString},
    commentstyle     = \color{codeComment}\itshape,
    numberstyle      = \tiny\color{codeNumber},
    numbers          = left,
    keepspaces       = true,
    columns          = fullflexible,
    upquote          = true,
    numbersep        = 8pt,
    frame            = single,
    rulecolor        = \color{codeFrame},
    framesep         = 6pt,
    xleftmargin      = 16pt,
    framexleftmargin = 16pt,
    breaklines       = true,
    breakatwhitespace= true,
    showstringspaces = false,
    tabsize          = 4,
    captionpos       = t,
    aboveskip        = 12pt,
    belowskip        = 12pt,
    literate         = {á}{{\'a}}1 {é}{{\'e}}1 {í}{{\'i}}1 {ó}{{\'o}}1
                       {ú}{{\'u}}1 {ã}{{\~a}}1 {õ}{{\~o}}1
                       {ç}{{\c{c}}}1 {Ç}{{\c{C}}}1
                       {ê}{{\^e}}1 {â}{{\^a}}1
}

% --- Listas ---
\usepackage{enumitem}
\setlist{
    topsep     = 0pt,
    itemsep    = 0pt,
    parsep     = 4pt,
    leftmargin = \parindent
}
\setlist[enumerate]{label=\arabic*.}
\setlist[itemize]{label=\textbullet}

% --- Referências e Hiperlinks ---
\usepackage[style=abnt, justify, maxbibnames=99, sorting=nty]{biblatex}
\setlength{\bibitemsep}{\baselineskip} % Espaçamento de 1 linha entre referências distintas (ABNT NBR 6023)
\addbibresource{referencias.bib}

\usepackage{xurl}
\urlstyle{same}
\usepackage{hyperref}
\hypersetup{
    colorlinks     = true,
    linkcolor      = black,
    citecolor      = black,
    filecolor      = black,
    urlcolor       = black,
    bookmarksdepth = 4
}

% --- Margens e Espaçamentos (ABNT) ---
\setlrmarginsandblock{3cm}{2cm}{*}
\setulmarginsandblock{3cm}{2cm}{*}
\checkandfixthelayout

\OnehalfSpacing
\setlength{\parindent}{1.25cm}
\setlength{\ABNTEXcitacaorecuo}{4cm}

\widowpenalty=10000
\clubpenalty=10000
\brokenpenalty=10000 % Impede palavras hifenizadas quebrando entre viradas de página

% --- Estilização Acadêmica Sóbria ---
\renewcommand{\ABNTEXchapterfont}{\normalfont\bfseries}
\renewcommand{\ABNTEXchapterfontsize}{\normalsize} % NBR 14724 exige fonte tamanho 12 uniforme
\renewcommand{\ABNTEXsectionfont}{\normalfont\bfseries}
\renewcommand{\ABNTEXsectionfontsize}{\normalsize}
\renewcommand{\ABNTEXsubsectionfont}{\normalfont\bfseries}
\renewcommand{\ABNTEXsubsectionfontsize}{\normalsize}
\setboolean{ABNTEXuppersection}{false} % Apenas capítulos ficam em caixa alta, seções secundárias em caixa baixa

% Nomes de elementos pré e pós-textuais em maiúsculas (ABNT NBR 14724)
\addto\captionsbrazilian{
    \renewcommand{\bibname}{REFERÊNCIAS}
    \renewcommand{\contentsname}{SUMÁRIO}
    \renewcommand{\listfigurename}{LISTA DE FIGURAS}
    \renewcommand{\listtablename}{LISTA DE TABELAS}
    \renewcommand{\listofquadrosname}{LISTA DE QUADROS}
    \renewcommand{\listadesiglasname}{LISTA DE ABREVIATURAS E SIGLAS}
}

\setlength{\beforechapskip}{18pt}
\setlength{\afterchapskip}{1.5\baselineskip}
\chapterstyle{abnt}

\renewcommand{\cftdot}{.}
\renewcommand{\cftdotsep}{3}
\setlength{\cftbeforechapterskip}{0.3em}
\setpnumwidth{2.5em} % Espaço garantido para numeração de páginas com 3 dígitos nas listas
\setrmarg{3.5em plus 1fil} % Alinhamento flexível nas listas (ToC/LoF/LoT) para evitar corte de página por Overfull \hbox

% ==================================================================
% MACROS DE IMPRESSÃO (CAPA E FOLHA DE ROSTO IFS)
% ==================================================================

\renewcommand{\imprimircapa}{%
  \begin{capa}%
    \center\bfseries % Todos os elementos da capa em negrito e centralizados (padrão institucional)
    \OnehalfSpacing
    
    \includegraphics[width=3.2cm]{logo_ifs.png}\par
    \vspace{0.5cm}
    
    {\normalsize\MakeUppercase{\imprimirinstituicao}\par}
    
    \vspace{3.0cm}
    
    {\normalsize\MakeUppercase{\imprimirautor}\par}
    
    \vspace{3.0cm}
    
    {\normalsize\MakeUppercase{\imprimirtitulo}\par}
    
    \vfill
    
    {\normalsize\MakeUppercase{\imprimirlocal}\par}
    {\normalsize\imprimirdata\par}
  \end{capa}
}

\makeatletter
% --- Variáveis Opcionais de Banca e Aprovação (preenchidas em src/metadados.tex) ---
\newcommand{\@dataaprovacao}{}
\newcommand{\dataaprovacao}[1]{\renewcommand{\@dataaprovacao}{#1}}
\newcommand{\@orientadorinstituicao}{Instituto Federal de Educação, Ciência e Tecnologia de Sergipe -- IFS}
\newcommand{\orientadorinstituicao}[1]{\renewcommand{\@orientadorinstituicao}{#1}}
\newcommand{\@coorientadorinstituicao}{}
\newcommand{\coorientadorinstituicao}[1]{\renewcommand{\@coorientadorinstituicao}{#1}}
\newcommand{\@bancaum}{}
\newcommand{\bancaum}[1]{\renewcommand{\@bancaum}{#1}}
\newcommand{\@bancauminstituicao}{}
\newcommand{\bancauminstituicao}[1]{\renewcommand{\@bancauminstituicao}{#1}}
\newcommand{\@bancadois}{}
\newcommand{\bancadois}[1]{\renewcommand{\@bancadois}{#1}}
\newcommand{\@bancadoisinstituicao}{}
\newcommand{\bancadoisinstituicao}[1]{\renewcommand{\@bancadoisinstituicao}{#1}}
\newcommand{\@bancatres}{}
\newcommand{\bancatres}[1]{\renewcommand{\@bancatres}{#1}}
\newcommand{\@bancatresinstituicao}{}
\newcommand{\bancatresinstituicao}[1]{\renewcommand{\@bancatresinstituicao}{#1}}

\renewcommand{\imprimirfolhaderosto}{%
  \begin{folhaderosto}
    \center
    \OnehalfSpacing
    
    \vspace*{2cm}
    
    {\normalsize\MakeUppercase{\imprimirautor}\par}
    
    \vspace{4.0cm}
    
    {\normalsize\bfseries\MakeUppercase{\imprimirtitulo}\par}
    
    \vspace{2.0cm}
    
    \begin{flushright}
    \begin{minipage}{8cm}
        \SingleSpacing
        \normalsize
        \imprimirpreambulo\par
        \vspace{0.5cm}
        Orientador: \imprimirorientador
        \ifdefempty{\imprimircoorientador}{}{%
            \par\vspace{0.2cm}
            Coorientador: \imprimircoorientador
        }
    \end{minipage}
    \end{flushright}
    
    \vfill
    
    {\normalsize\MakeUppercase{\imprimirlocal}\par}
    {\normalsize\imprimirdata\par}
  \end{folhaderosto}
}

% --- Folha de Aprovação Institucional (IFS / ABNT NBR 14724) ---
\newcommand{\imprimirfolhadeaprovacao}{%
  \begin{folhadeaprovacao}
    \begin{center}
      {\normalsize\MakeUppercase{\imprimirautor}\par}
      \vspace{1.2cm}
      {\normalsize\bfseries\MakeUppercase{\imprimirtitulo}\par}
    \end{center}
    \vspace{1.0cm}
    \begin{flushright}
      \begin{minipage}{8cm}
        \SingleSpacing
        \normalsize
        \imprimirpreambulo
      \end{minipage}
    \end{flushright}
    \vspace{0.8cm}
    \noindent Aprovado em: \ifdefempty{\@dataaprovacao}{\underline{\hspace{1.2cm}} de \underline{\hspace{3.2cm}} de \imprimirdata.}{\@dataaprovacao}\par
    \vspace{0.8cm}
    \begin{center}
      {\normalsize\bfseries BANCA EXAMINADORA}
    \end{center}
    \setlength{\ABNTEXsignwidth}{11cm}
    \setlength{\ABNTEXsignthickness}{0.8pt}
    \setlength{\ABNTEXsignskip}{0.9cm}
    \assinatura{\textbf{\imprimirorientador} (Orientador) \\ \@orientadorinstituicao}
    \ifdefempty{\imprimircoorientador}{}{%
      \assinatura{\textbf{\imprimircoorientador} (Coorientador) \\ \@coorientadorinstituicao}%
    }
    \ifdefempty{\@bancaum}{}{%
      \assinatura{\textbf{\@bancaum} (Avaliador) \\ \@bancauminstituicao}%
    }
    \ifdefempty{\@bancadois}{}{%
      \assinatura{\textbf{\@bancadois} (Avaliador) \\ \@bancadoisinstituicao}%
    }
    \ifdefempty{\@bancatres}{}{%
      \assinatura{\textbf{\@bancatres} (Avaliador) \\ \@bancatresinstituicao}%
    }
    \vfill
    \begin{center}
      {\normalsize\MakeUppercase{\imprimirlocal}\par}
      {\normalsize\imprimirdata\par}
    \end{center}
  \end{folhadeaprovacao}
}

% --- Ficha Catalográfica (PDF fornecido pela biblioteca do IFS) ---
\newcommand{\inserirfichacatalografica}[1]{%
  \IfFileExists{#1}{%
    \begin{fichacatalografica}
      \includepdf[pages=1]{#1}
    \end{fichacatalografica}
  }{}%
}
\makeatother

% ==================================================================
% FIM DAS CONFIGURAÇÕES
% ==================================================================
